\documentclass[11pt]{article}
\usepackage{amsmath, amssymb}
\usepackage[margin=1in]{geometry}
\usepackage{fancyhdr}
\usepackage{hyperref}
\usepackage{array}

\hypersetup{
    colorlinks=true,
    linkcolor=blue,
    urlcolor=blue,
    citecolor=green
}

% ---- Fill these in ----
\newcommand{\courseName}{NEU 314}
\newcommand{\taName}{Keith Murray}
\newcommand{\sectionInfo}{Precept B01}
% ------------------------

\newcommand{\E}{\mathbb{E}}

\title{Precept 6: Probability Exercises}
\date{\today}

\pagestyle{fancy}
\fancyhf{}
\lhead{\courseName}
\chead{\sectionInfo}
\rhead{\today}
\cfoot{\thepage}

\begin{document}

\maketitle
\thispagestyle{fancy}

\section*{Problem}
The table below gives $P(x,y)$.
\begin{center}
\renewcommand{\arraystretch}{1.6}
\begin{tabular}{cc|>{\centering\arraybackslash}p{1.6cm}|>{\centering\arraybackslash}p{1.6cm}|>{\centering\arraybackslash}p{1.6cm}|}
    \multicolumn{2}{c}{} & \multicolumn{3}{c}{$x$} \\
    \multicolumn{2}{c}{} & \multicolumn{1}{c}{1} & \multicolumn{1}{c}{2} & \multicolumn{1}{c}{3} \\
    \cline{3-5}
    & 1 & 0.02 & 0.06 & 0.04 \\
    \cline{3-5}
    $y$ & 2 & 0.10 & 0.36 & 0.22 \\
    \cline{3-5}
    & 3 & 0.04 & 0.10 & 0.06 \\
    \cline{3-5}
\end{tabular}
\end{center}

\begin{enumerate}
    \item Verify this is a proper joint distribution.
    \item What is the marginal $P(x)$?
    \item What is the marginal $P(y)$?
    \item What is the conditional $P(y \mid x=1)$?
    \item What is the conditional $P(x \mid y<3)$?
    \item Compute $\displaystyle \E[x]=\sum_{x=1}^{3} x\,P(x)$. What is this?
\end{enumerate}

\newpage
\section*{Solution}

\subsection*{Question 1}
To get a proper joint distribution, we need two things:
\begin{enumerate}
    \item All entries must be non-negative,
    \item All entries must sum to 1.
\end{enumerate}
Indeed, all entries are non-negative. To show they sum to one, we can write
\[
\underbrace{(0.02+0.06+0.04)}_{0.12}+\underbrace{(0.10+0.36+0.22)}_{0.68}+\underbrace{(0.04+0.10+0.06)}_{0.20}=1
\]

\subsection*{Question 2}
To get the marginal $P(x)$, we sum out $y$, meaning that we sum down each column.
\[
P(x=1)=0.02+0.10+0.04=0.16
\]
\[
P(x=2)=0.06+0.36+0.10=0.52
\]
\[
P(x=3)=0.04+0.22+0.06=0.32
\]
As a sanity check, $0.16+0.52+0.32=1$.

\subsection*{Question 3}
To get the marginal $P(y)$, we sum out $x$ by summing across each row.
\[
P(y=1)=0.02+0.06+0.04=0.12
\]
\[
P(x=2)=0.10+0.36+0.22=0.68
\]
\[
P(x=3)=0.04+0.10+0.06=0.20
\]
As a sanity check, $0.12+0.68+0.20=1$.

\subsection*{Question 4}
The conditional distribution is
\[
P(y\mid x=1)=\frac{P(x=1,y)}{P(x=1)}
\]
Using $P(x=1)=0.16$, we can write
\[
P(y=1\mid x=1)=\frac{0.02}{0.16}=0.125,\quad
P(y=2\mid x=1)=\frac{0.10}{0.16}=0.625,\quad
P(y=3\mid x=1)=\frac{0.04}{0.16}=0.25
\]
which sums to 1, as every distribution should.

\subsection*{Question 5}
The conditional distribution is
\[
P(x\mid y<3)=\frac{P(x,y<3)}{P(y<3)}
\]
First, let's compute $P(y< 3)$. This is just summing from the marginal distribution
\[
P(y< 3) = P(y=1) + P(y=2) = 0.12 + 0.68 = 0.80
\]
Now, we can compute the following 
\begin{gather*}
    P(x=1\mid y<3)=\frac{0.02+0.10}{0.80}=0.15,\\
    P(x=2\mid y<3)=\frac{0.06+0.36}{0.80}=0.525,\\
    P(x=3\mid y<3)=\frac{0.04+0.22}{0.80}=0.325
\end{gather*}

\subsection*{Question 6}
\[
\E[x]=1(0.16)+2(0.52)+3(0.32)=0.16+1.04+0.96=2.16
\]
This is the mean!


\end{document}